\documentclass[a4paper,11pt,oneside]{report}

\newif\ifrelease % new boolean variable release. True = include some fancy content
\releasetrue % and set it

\usepackage[english]{babel}
\usepackage[utf8]{inputenc}
\usepackage[T1]{fontenc}
\usepackage{amsmath}
\usepackage{epsfig}
\usepackage{graphicx}
\ifrelease
	\usepackage{pdfpages}
\fi
\usepackage[pagebackref=true]{hyperref} % tento balicek by mel byt na konci baliku!

\hypersetup{
	pdfauthor={Matej Laitl},
	pdftitle={Implementation environment for Bayesian filtering algorithms}
}

%% Nastavení zrcadla sazby
\usepackage{calc}
\setlength{\textheight}{9in}
\setlength{\textwidth}{6in}
\setlength\oddsidemargin{(\paperwidth-\textwidth)/2 - 1in}
\setlength\evensidemargin{(\paperwidth-\textwidth)/2 - 1in}
\setlength\topmargin{(\paperheight-\textheight-\headheight-\headsep-\footskip)/2 - 1in}

% TODO: je toto potreba?
%\parindent=0pt % odsazení 1. řádku odstavce
%\parskip=7pt   % mezera mezi odstavci

\ifrelease
	\hypersetup{pdfborder={0 0 0}} % no borders around links
\else
	\hypersetup{colorlinks=true} % colour links instead of borders
\fi

\begin{document}

\ifrelease
	\newcommand{\cvut}{České vysoké učení technické v~Praze}
\newcommand{\fjfi}{Fakulta jaderná a fyzikálně inženýrská}
\newcommand{\km}{Katedra matematiky}
\newcommand{\obor}{Inženýrská Informatika}
\newcommand{\zamereni}{Softwarové inženýrství a matematická informatika}

\newcommand{\nazevcz}{Prostředí pro implementaci algoritmů Bayesovské filtrace}
\newcommand{\nazeven}{Implementation environment for Bayesian filtering algorithms}
\newcommand{\autor}{Matěj Laitl}
\newcommand{\rok}{2011}
\newcommand{\vedouci}{Ing. Václav Šmídl, Ph.D.}

\newcommand{\pracovisteVed}{Oddělení adaptivních systémů \\
	Ústav teorie informace a automatizace \\
	Akademie věd České republiky}
\newcommand{\konzultant}{---}
\newcommand{\pracovisteKonz}{}

\newcommand{\klicova}{Bayesovská filtrace, softwarová analýza, Python, Cython} % TODO: dalsi?
\newcommand{\keyword}{Bayesian filtering, software analysis, Python, Cython}
% Abstrakt práce: (cca 7 vět, min. 80 slov)
\newcommand{\abstrCZ}{Bayesovská filtrace je velmi použitelný přístup k odhadování dynamických
systémů, který má mnoho aplikací v robotice, environmentálních simulacích a dalších oborech. Tato
práce předkládá popis Kalmanova filtru, částicového filtru a marginalizovaného částicového filtru,
který využívá myšlenky obou předchozích. Poté je provedena softwarová analýza, která má za úkol
určit nejvhodnější implementační prostředí pro softwarovou knihovnu poskytující metody Bayesovské
filtrace. Nejprve jsou posouzeny obecné přístupy k vývoji software, z nichž je jako nejpoužitelnější
zvolen objektově orientovaný, následuje porovnání programovacích jazyků C++, MATLAB a Python.
Python zvítězí a v je něm napsána knihovna pro Bayesovskou filtraci nazvaná PyBayes; nakonec je
dosaženo výrazného zrychlení Pythonového kódu použitím Cythonu.}
\newcommand{\abstrEN}{Bayesian filtering is a vital approach to dynamic system estimation that has
wide applications in robotics, environmental simulations and much more. This text briefly introduces
the Kalman filter, the particle filter and the marginalized particle filter which combines ideas of
both. Software analysis is performed with the aim to identify the most suitable implementation
environment for a software library for Bayesian filtering. Various paradigms of software development
and are discussed among which the object-oriented programming is chosen; C++, MATLAB and Python
languages are evaluated. Python is determined most appropriate and a Python software library named
PyBayes is developed; dramatic performance gains are measured when Cython is used to speed up
Python.}

%%% zde zacina kresleni dokumentu

% titulní strana
\thispagestyle{empty}

\begin{center}
	{\Large  \bf  \cvut\\[2mm] \fjfi }
	\vspace{10mm}

	\begin{tabular}{c}
	{\bf \km}\\
	{\bf Obor: \obor}\\
	{\bf Zaměření: \zamereni}
	\end{tabular}

	\vspace{10mm} \epsfysize=20mm  \epsffile{cvut-logo-bw-600} \vspace{15mm}

	{\LARGE
	\textbf{\nazevcz}
	\par}

	\vspace{5mm}

	{\LARGE
	\textbf{\nazeven}
	\par}

	\vspace{30mm}
	{\Large BAKALÁŘSKÁ PRÁCE}

\end{center}

\vfill
{\large
\begin{tabular}{rl}
Vypracoval: & \autor\\
Vedoucí práce: & \vedouci\\
Rok: & \rok
\end{tabular}
}

% zadání bakalářské práce
\newpage
\thispagestyle{empty} Před svázáním místo téhle stránky \fbox{vložíte zadání práce} s podpisem
děkana (bude to jediný oboustranný list ve Vaší práci) !!!!

% prohlášení
\newpage
\thispagestyle{empty}
~
\vfill


{\bf Prohlášení}

\vspace{0.5cm}
Prohlašuji, že jsem svou bakalářskou práci vypracoval samostatně a použil jsem pouze podklady
(literaturu, projekty, SW atd.) uvedené v přiloženém seznamu.

\vspace{5mm}V Praze dne ....................\hfill
    \begin{tabular}{c}
    ........................................\\
    \autor
    \end{tabular}

% poděkování
\newpage
\thispagestyle{empty}
~
\vfill

{\bf Poděkování}

\vspace{5mm}
Děkuji Ing. Václavu Šmídlovi, Ph.D. za vedení mé bakalářské práce, velmi prozíravě položené
softwarové otázky a za poukázání na softwarové projekty (Cython, Sphinx, \ldots), které se ukázaly
být klíčové pro moji bakalářskou práci a softwarový projekt.

\begin{flushright}
\autor
\end{flushright}

% strana s abstraktem
\newpage
\thispagestyle{empty}

\newbox\odstavecbox
\newlength\vyskaodstavce
\newcommand\odstavec[2]{%
    \setbox\odstavecbox=\hbox{%
         \parbox[t]{#1}{#2\vrule width 0pt depth 4pt}}%
    \global\vyskaodstavce=\dp\odstavecbox
    \box\odstavecbox}
\newcommand{\delka}{120mm}

\noindent\begin{tabular}{@{}ll}
  {\em Název práce:} & ~ \\
  \multicolumn{2}{@{}l}{\odstavec{\textwidth}{\bf \nazevcz}} \\[5mm]
  {\em Autor:} & \autor \\[5mm]
  {\em Obor:} & \obor \\
  {\em Druh práce:} & Bakalářská práce \\[5mm]
  {\em Vedoucí práce:} & \odstavec{\delka}{\vedouci \\ \pracovisteVed} \\[5mm]
  {\em Konzultant:} & \odstavec{\delka}{\konzultant \\ \pracovisteKonz} \\[5mm]
  \multicolumn{2}{@{}l}{\odstavec{\textwidth}{{\em Abstrakt:} ~ \abstrCZ \\ }} \\[5mm]
  {\em Klíčová slova:} & \odstavec{\delka}{\klicova} \\[10mm]

  {\em Title:} & ~\\
  \multicolumn{2}{@{}l}{\odstavec{\textwidth}{\bf \nazeven}}\\[5mm]
  {\em Author:} & \autor \\[5mm]
  \multicolumn{2}{@{}l}{\odstavec{\textwidth}{{\em Abstract:} ~ \abstrEN \\ }} \\[5mm]
  {\em Key words:} & \odstavec{\delka}{\keyword}
\end{tabular}
 % include some fancy start pages
\fi


% obsah
\newpage
\tableofcontents


\chapter*{Introduction} \addcontentsline{toc}{chapter}{Introduction}

TODO motivatin for bayes filtration + a need for a convenient library (rapid prototyping vs. speed)

Decision-making being a logical and natural ``next step''...


\chapter{Basics of Bayesian estimation}

\section{Kalman Filter}

\section{Particle Filter}

\section{Marginalized Particle Filter}


\chapter{Software analysis}

\section{Requirements}

Ideal library for Bayesian filtering would posses following properties...:

\section{Programming paradigms}

Interpreted vs. Compiled

Object-oriented, procedural and Functional

pass-by reference vs. copy-on-write (Matlab)

\section{C++}

BDM

 - advantages of C (speed, C prevalence (many optimised libraries, BLAS, LAPACK.., OpenMP)

 - disadvantages of C in our ``situation'' (steep learning curve, coplexity because of low-levelness
   high initial barriers (need to have compiler, libraries...), inconveniently long edit/build/test
   process)

\section{Matlab}

BDM (partially?)

 - advantages (popularity, existing toolboxes, rapid development (high-level)

 - disadv: strict copy-on-write, problematic object model (not in original design), difficulties
           interfacing existing C (F) code

\section{Python}

NumPy.... parallelisation (approaches, improvements in Py 3.2) - GIL.. Py3k

\section{Cython}

general info etc... extension types, building, ease of interfacing C (and F) code, .pxd files,
NumPy support

\subsection{Gradual Optimisation and Parallelisation}

how can optimisaion be approached (gradually) and why this approach is superior

integrate\_python\_cython example (``100x'' speedup for a special (very simple) case)

\subsection{Parallelisation}

integrate\_python\_cython patched with OpenMP (13x speedup in 16-core system)

prange CEP

\subsection{Pure Python mode}

About it and why it should be used in a hypothetical bayesian python library

\subsection{Limitations}

2 types:

	not-supported code (few, but bad, ongoing work)

	not-optimised code (much more work needed, but not hard to fix in most cases)

		- exception handling (functions returning void etc)

		- limitations of pure python mode in regards to traditional .pyx files


\chapter{Finding a Way to a Suitable Bayesian Library}

TODO: maybe not a chapter, but a section in another one (also better name for chapter)

\section{The choice of Python/Cython}

\section{Existing Python Libraries for Bayesian estimation/decision making}

Survey of existing libraries ... not suitable

The need to implement a new one :-)


\chapter{The PyBayes Library}

Introduction, general directions, future considerations

+ open development on github, open-source

\section{Interpreted and Compiled}

\section{Library Layout}

\subsection{Random Variable Meta-representation}

Why it is needed (ref to ProdCPdf)

\subsection{Probability Density Functions}

Nice UML diagrams! (better more smaller UMLs than one big) One for general pdf layut, one for
AbstractGaussPdf family, one for AbstractEmpPdf family

\subsection{Bayesian Filters}

Nice graph of a run of a particle filter (Mirda has the plotting code)

similar of marginalized particle filter? (gausses would be plotted vertically)

\section{Documentation, Testing and Profiling}

TODO: move above Library Layout?

Documenting using Sphinx, approach to documentation (mathematician-oriented), math in documentation

Testing - the separation of

	- tests: test one class in isolation, quick, determinism (would be good, not achievable)

	- stresses: test a great portion of code at once, run longer, non-determinism..

Profiling python/cython - how, existing support in PyBayes

	- how to correct profiling-induced overhead

\section{Comparison with BDM}


\chapter*{Conclusion} \addcontentsline{toc}{chapter}{Conclusion}


% použitá literatura
\nocite{*} % TODO: jen docasne pro to, aby se zobrazily vsechny veci z BibTex databaze

\clearpage % so that the contents link mentions correct page
\phantomsection % so that hyperref makes correct reference
\addcontentsline{toc}{chapter}{\bibname}
\bibliographystyle{plain}
\bibliography{bibliography}

\ifrelease
	% přílohy
	\appendix % aby LaTeX cisloval jinak
	\clearpage % so that the contents link mentions correct page
	\phantomsection % so that hyperref makes correct reference
	\addcontentsline{toc}{chapter}{\appendixname}

	\part*{\appendixname}

	\includepdf[pages=-]{../doc/_build/latex/PyBayes.pdf}
\fi

\end{document}
